\chapter{Понятие власти}

\section{Тезис: власть как логика формы}

Власть — это не функция силы, а форма, определяющая, каким образом сила становится действием. Она не находится вне субъекта и не привносится извне, но проявляется как внутренняя структура воли, формализованная в отношении между я и другим.

Она есть форма актуализации различия между тем, кто утверждает, и тем, кто подчиняется. Не команда делает властью, а форма команды, получившая признание как необходимая.

Власть первична по отношению к содержанию, ибо она задаёт само условие возможности действия: не «что» осуществляется, а «кто» может осуществить.

\section{Антитезис: власть как инструмент}

Если в логике формы власть предстает как структурное условие действия, то в логике средства она редуцируется до функции, до инструмента воли. Здесь субъект уже не соотносит себя с Другим в рамках признания, а стремится к достижению цели путём доминирования, манипуляции, управления. Власть становится продолжением желания, лишённого необходимости быть признанным.

В этом подходе власть — это не форма, а средство воздействия: «инструмент принуждения», «ресурс подчинения», «рычаг контроля». Она становится орудием, подчинённым цели, как нож подчинён намерению резать. Этим её понятие утрачивает самодвижение и становится механическим.

Такое понимание власти характерно для утилитарных и прагматических теорий: у Макиавелли, у Гоббса, у современных технократов. Но тем самым власть теряет онтологическую глубину и превращается в технологию — отчуждённую, передаваемую, измеримую. Она более не проявление воли, а набор приёмов и процедур.

Такое видение власти делает её количественной, а не качественной: власть — это «больше» или «меньше» контроля, влияния, ресурсов. Субъект здесь исчезает как носитель формы — остаётся лишь оператор.



\section{Синтез: власть как процесс становления}
% Здесь пока оставим пусто
